\clearpage
\fancyhead[L]{\small JUNO.}
\hypertarget{native-263-lower}{}
\begin{center}
JUNO.

\rule{0.6\linewidth}{0.4pt}

Monatliche Correspondenz. Band X. S. 552–555. 1804 December.

\rule{0.6\linewidth}{0.4pt}
\end{center}

\noindent—\quad—\quad—\par

Leider hinderte uns der seit dem 6. November hier stets umwölkte Himmel, die Juno fortdauernd
zu beobachten und dem Dr. \textsc{Gauss} neue Data zur Rectification seiner bereits berechneten Elemente zu
liefern; allein zwei von uns am 20. und 21. Oct. gemachte Beobachtungen, die wir ihm mittheilten, und
die $1\frac{1}{2}$ Minute von dessen II. Elementen abwichen, waren diesem eben so unermüdeten als scharfsinnigen Astronomen hinlängliche Veranlassung, sogleich neue III. Elemente für die Juno zu berechnen, die
er uns mit folgenden Bemerkungen begleitet überschickte: „Seit meinem letzten Briefe habe ich, mit
Hülfe Ihrer mir gütigst mitgetheilten Beobachtungen der Juno vom 20. und 21. Oct., die von den zweiten Elementen bereits $1\frac{1}{2}$ Minute differirten, folgende neue III. Elemente berechnet:

\begin{center}
\begingroup\fontsize{8}{10}\selectfont\setlength{\tabcolsep}{3pt}\renewcommand{\arraystretch}{1.1}
\begin{tabular}{lr}
Epoche 1804 Sept. 30 $0^{\mathrm u}$ im Meridian von Seeberg . . & $22^\circ\ 34'\ 48''$ \\
tägliche mittlere Bewegung . . . . . . . . . . . & $842''75$ \\
Sonnenferne . . . . . . . . . . . . . . . . & $233^\circ\ 56'\ 6''$ \\
Logarithmus der halben grossen Axe . . . . . . . & $0.426699$ \\
Excentricität . . . . . . . . . . . . . . . . & $0.263182$ \\
aufsteigender Knoten . . . . . . . . . . . . . & $171^\circ\ 0'\ 0''$ \\
Neigung der Bahn . . . . . . . . . . . . . . & $12\ 52\ 48$ \\
\end{tabular}\endgroup
\end{center}

„Ich habe diese Elemente mit Ihren sämmtlichen Beobachtungen verglichen, und folgende Übereinstimmung gefunden:
\SourceLeaf{264}
\begin{center}
\begingroup\fontsize{8}{10}\selectfont\setlength{\tabcolsep}{3pt}\renewcommand{\arraystretch}{1.1}
\begin{tabular}{r|r|r|r|r|r}
1804 & \shortstack{Mittlere Zeit\\auf\\Seeberg} & \shortstack{Berechnete gerade\\Aufsteigung\\der $\JunoSymbol$} & \shortstack{Berechn. südl.\\Abweichung\\der $\JunoSymbol$} & \multicolumn{2}{c}{\shortstack{Unterschied\\R—B\\in AR.\qquad in Decl.}}\\ \hline
Sept. 13 & $12^{\mathrm u}\ 31^{\mathrm m}\ 59^{\mathrm s}$ & $0^\circ\ 44'\ 55''8$ & $1^\circ\ 52'\ 34''3$ & $-0''6$ & $-2''7$ \\
14 & $12\ 27\ 27$ & $\phantom{000}\ 35\ 43.2$ & $2\ 5\ 35.3$ & $+2.2$ & $-0.2$ \\
15 & $12\ 22\ 53$ & $\phantom{000}\ 26\ 20.8$ & $\phantom{0}\ 18\ 42.2$ & $+0.8$ & $+3.7$ \\
17 & $12\ 13\ 45$ & $\phantom{000}\ 7\ 10.6$ & $\phantom{0}\ 45\ 9.2$ & $+2.5$ & $-3.9$ \\
18 & $12\ 9\ 10$ & $359\ 57\ 24.5$ & $\phantom{0}\ 58\ 28.2$ & $+4.0$ & $-2.7$ \\
20 & $.\ .\ .\ .$ & $\phantom{000}\ 37\ 35.7$ & $3\ 25\ 32.7$ & $.\ .\ .\ .$ & $+9.2$ \\
23 & $11\ 46\ 11$ & $\phantom{000}\ 7\ 25.9$ & $4\ 5\ 22.9$ & $+11.0$ & $-11.6$ \\
27 & $11\ 27\ 46$ & $358\ 27\ 7.7$ & $\phantom{0}\ 58\ 29.8$ & $+7.0$ & $-27.4$ \\
28 & $11\ 23\ 11$ & $\phantom{000}\ 17\ 9.6$ & $5\ 11\ 35.1$ & $+3.2$ & $-8.8$ \\
30 & $11\ 14\ 1$ & $357\ 57\ 28.6$ & $\phantom{0}\ 37\ 28.2$ & $-1.9$ & $-15.5$ \\
Oct. 2 & $11\ 4\ 52$ & $\phantom{000}\ 38\ 15.5$ & $6\ 2\ 50.8$ & $+0.5$ & $-8.7$ \\
4 & $10\ 55\ 46$ & $\phantom{000}\ 19\ 40.9$ & $\phantom{0}\ 27\ 36.7$ & $+1.4$ & $-2.0$ \\
5 & $10\ 51\ 15$ & $\phantom{000}\ 10\ 41.1$ & $\phantom{0}\ 39\ 43.8$ & $-0.4$ & $+3.9$ \\
6 & $10\ 46\ 44$ & $\phantom{000}\ 1\ 54.0$ & $\phantom{0}\ 51\ 39.0$ & $-2.3$ & $+6.8$ \\
10 & $10\ 28\ 50$ & $356\ 29\ 14.9$ & $7\ 37\ 10.5$ & $+0.8$ & $+19.4$ \\
12 & $10\ 20\ 0$ & $\phantom{000}\ 14\ 37.8$ & $7\ 58\ 30.8$ & $-5.7$ & $+4.3$ \\
20 & $9\ 45\ 34$ & $355\ 29\ 55.3$ & $9\ 12\ 58.8$ & $+3.4$ & $-4.6$ \\
21 & $9\ 41\ 23$ & $\phantom{000}\ 26\ 2.7$ & $\phantom{0}\ 20\ 45.8$ & $+2.6$ & $+2.6$ \\
\end{tabular}\endgroup
\end{center}

„Dr. \textsc{Maskelyne} hat mir noch folgende drei Beobachtungen mitzutheilen die Güte gehabt:

\begin{center}
\begingroup\fontsize{9}{11}\selectfont\setlength{\tabcolsep}{3pt}\renewcommand{\arraystretch}{1.1}
\begin{tabular}{r|r|r|r}
1804 & \shortstack{Mittlere Zeit\\in\\Greenwich} & \shortstack{Scheinbare gerade\\Aufsteigung\\der $\JunoSymbol$} & \shortstack{Scheinbare süd-\\liche Abwei-\\chung der $\JunoSymbol$}\\ \hline
Oct. 5 & $10^{\mathrm u}\ 51^{\mathrm m}\ 6^{\mathrm s}$ & $357^\circ\ 10'\ 24''6$ & $6^\circ\ 40'\ 1''1$ \\
9 & $10\ 33\ 9$ & $356\ 36\ 46.9$ & $7\ 26\ 21.3$ \\
17 & $9\ 58\ 10$ & $355\ 43\ 45.6$ & $8\ 47\ 18.5$ \\
\end{tabular}\endgroup
\end{center}

„Damit und den frühern Beobachtungen stimmen die III. Elemente, wie folgt:

\begin{center}
\begingroup\fontsize{8}{10}\selectfont\setlength{\tabcolsep}{3pt}\renewcommand{\arraystretch}{1.1}
\begin{tabular}{r|r|r|r|r|r}
1804 & \shortstack{Mittlere Zeit\\in\\Greenwich} & \shortstack{Berechn. gerade\\Aufsteigung\\der $\JunoSymbol$} & \shortstack{Berechn. südl.\\Abweichung\\der $\JunoSymbol$} & \multicolumn{2}{c}{\shortstack{Unterschied\\R—B\\in AR.\qquad in Decl.}}\\ \hline
Sept. 25 & $11^{\mathrm u}\ 36^{\mathrm m}\ 50^{\mathrm s}$ & $358^\circ\ 46'\ 55''7$ & $4^\circ\ 32'\ 26''8$ & $+7''7$ & $-1''4$ \\
29 & $11\ 18\ 27$ & $358\ 6\ 58.7$ & $5\ 24\ 58.3$ & $+6.4$ & $+1.6$ \\
Oct. 5 & $10\ 51\ 6$ & $357\ 10\ 24.7$ & $6\ 40\ 6.0$ & $+0.1$ & $+4.9$ \\
9 & $10\ 33\ 9$ & $356\ 36\ 46.4$ & $7\ 26\ 28.0$ & $-0.5$ & $+6.7$ \\
17 & $9\ 58\ 10$ & $355\ 43\ 46.0$ & $8\ 47\ 27.1$ & $+0.4$ & $+8.6$ \\
\end{tabular}\endgroup
\end{center}

„Nach den III. Elementen steht der künftige Lauf der \emph{Juno} folgendermassen: [von 1804 Nov. 2 bis
1805 Jan. 4 mit dreitägigen Zwischenzeiten, gerade Aufsteigung, Abweichung und Lichtstärke.]

„Bei der Lichtstärke ist diejenige zur Einheit angenommen worden, die der Planet in der Distanz 1 von der Sonne und Erde haben würde. Nach demselben Maassstabe war sie

\begin{center}
\begingroup\fontsize{9}{11}\selectfont\setlength{\tabcolsep}{3pt}\renewcommand{\arraystretch}{1.1}
\begin{tabular}{lrl}
den 5. Sept. & $0.1378$ &  \\
— 12. — & $0.1484$ &  \\
— 3. Oct. & $0.1640$ & am grössten in diesem Jahre \\
— 21. — & $0.1543$ &  \\
\end{tabular}\endgroup
\end{center}

„Es wird interessant sein, zu sehen, wie lange \emph{Juno} diesmal sichtbar bleiben wird. In der nächsten künftigen Opposition im Anfange März 1806, im Sternbilde des Löwen, erreicht sie nur ein Viertel von der grössten Helligkeit dieses Jahres. Für den 31. December 1805 finde ich ihren Ort
\SourceLeaf{265}
\begin{center}
AR $176^\circ\ 45'$\quad Decl. südl. $2^\circ\ 44'$\quad Lichtstärke $0.0284$.
\end{center}

\Needspace{17\baselineskip}
\indent Sämmtliche Beobachtungen, die wir zu erhalten im Stande waren, sind folgende

\begin{center}
\begingroup\fontsize{9}{11}\selectfont\setlength{\tabcolsep}{3pt}\renewcommand{\arraystretch}{1.1}
\begin{tabular}{r|r|r|r}
1804 & \shortstack{Mittlere Zeit\\auf\\Seeberg} & \shortstack{Scheinbare\\gerade Aufsteig.\\der $\JunoSymbol$} & \shortstack{Scheinbare\\südl. Abweich.\\der $\JunoSymbol$}\\ \hline
Oct. 23 & $9^{\mathrm u}\ 33^{\mathrm m}\ 4^{\mathrm s}338$ & $355^\circ\ 19'\ 23''69$ & $9^\circ\ 35'\ 53''6$ \\
24 & $9\ 28\ 57.709$ & $355\ 16\ 42.51$ & $.\ .\ .\ .$ \\
30 & $9\ 4\ 54.644$ & $355\ 9\ 47.25$ & $10\ 18\ 43.8$ \\
Nov. 5 & $8\ 41\ 56.099$ & $355\ 19\ 2.57$ & $.\ .\ .\ .$ \\
18 & $7\ 55\ 51.375$ & $356\ 34\ 46.58$ & $10\ 58\ 16.2$ \\
\end{tabular}\endgroup
\end{center}

\hfill[\textsc{von Zach.}]
